\chapter{Topolojik Uzaylar}
\label{ch:topological_spaces}

%% ============================================================
\section{Konu}
%% ============================================================

\subsection{Sezgisel Giriş}

Topoloji, bir kümedeki noktaların birbirine ``yakın'' olup olmadığını, aralarındaki
sürekliliği ve şekilsel özellikleri, mesafe kavramına gerek duymadan inceleyen matematiksel
bir disiplindir. Bunu mümkün kılan temel araç \emph{açık küme} kavramıdır.

Düşünelim: gerçek doğruda $(a, b)$ aralığı ``açık''tır; her noktasının çevresinde küçük bir
açık aralık daha bulunur. Bu sezgiyi soyutlayarak \emph{herhangi} bir küme üzerinde topoloji
tanımlamak mümkündür.

\subsection{Formal Tanım}

\begin{definition}[Topolojik Uzay]
\label{def:topological_space}
Bir $X$ kümesi ve $\tau \subseteq \mathcal{P}(X)$ ailesi verilsin.
$(X, \tau)$ bir \textbf{topolojik uzay}, $\tau$ ise $X$ üzerinde bir \textbf{topoloji}
olarak adlandırılır, eğer aşağıdaki üç aksiyom sağlanıyorsa:
\begin{enumerate}
  \item[(T1)] $\emptyset \in \tau$ ve $X \in \tau$
  \item[(T2)] $U, V \in \tau \Rightarrow U \cap V \in \tau$
  \item[(T3)] $\{U_\alpha\}_{\alpha \in I} \subseteq \tau \Rightarrow
              \displaystyle\bigcup_{\alpha \in I} U_\alpha \in \tau$
\end{enumerate}
$\tau$'nun elemanları \textbf{açık küme} olarak adlandırılır.
\end{definition}

\subsection{Temel Örnekler}

\begin{table}[h]
\centering
\begin{tabular}{lll}
\toprule
\textbf{Topoloji} & \textbf{Tanım} & \textbf{Açık kümeler} \\
\midrule
Ayrık (discrete) & Her alt küme açık & $\tau = \mathcal{P}(X)$ \\
İndirgenmiş (indiscrete) & Yalnızca $\emptyset$ ve $X$ açık & $\tau = \{\emptyset, X\}$ \\
Kosonlu (cofinite) & $U$ açık $\iff$ $U=\emptyset$ veya $X\setminus U$ sonlu & — \\
Sierpiński & $X=\{0,1\}$, $\tau = \{\emptyset, \{1\}, X\}$ & 3 açık küme \\
\bottomrule
\end{tabular}
\caption{Standart topoloji örnekleri}
\end{table}

\subsection{Baz ve Alt-Baz}

\begin{definition}[Baz]
\label{def:basis}
$\mathcal{B} \subseteq \tau$ ailesi, her $U \in \tau$ ve her $x \in U$ için
$x \in B \subseteq U$ sağlayan bir $B \in \mathcal{B}$ varsa $\tau$'ya \textbf{baz} denir.
\end{definition}

\noindent\textbf{Baz Koşulları:}
\begin{enumerate}
  \item[(B1)] $\bigcup \mathcal{B} = X$
  \item[(B2)] $B_1, B_2 \in \mathcal{B}$ ve $x \in B_1 \cap B_2$ ise bir $B_3 \in \mathcal{B}$
              vardır: $x \in B_3 \subseteq B_1 \cap B_2$
\end{enumerate}

\begin{definition}[Alt-Baz]
$\mathcal{S} \subseteq \mathcal{P}(X)$ ailesi; sonlu kesişimleri baz, onların birleşimleri
ise $\tau$'yu oluşturur.
\end{definition}

\subsection{Sonlu ve Sonsuz Uzaylar}

\texttt{pytop}'ta iki temel sınıf bulunur:
\begin{itemize}
  \item \textbf{\texttt{FiniteTopologicalSpace}:} Taşıyıcı $(X)$ sonlu; topoloji açıkça
        \texttt{frozenset}'ler koleksiyonu olarak saklanır. Tüm hesaplamalar tam ve kesindir.
  \item \textbf{Sonsuz türevler:} Taşıyıcı simgesel (\texttt{'R'}, \texttt{'N'});
        \texttt{topology=None}; etiket tabanlı sembolik çıkarım yapılır.
\end{itemize}

%% ============================================================
\section{Teoremler}
%% ============================================================

\begin{theorem}[Aksiyomların Yeterliliği]
\label{thm:axiom_sufficiency}
$(X, \tau)$ bir topolojik uzay olsun. T1--T3 aksiyomları, sonlu herhangi bir açık küme
ailesinin kesişiminin $\tau$'da yer aldığını garanti eder:
\[
  U_1, \ldots, U_n \in \tau \Rightarrow \bigcap_{i=1}^{n} U_i \in \tau.
\]
\textit{Not:} Sonsuz kesişim gerekmez. Gerçek doğruda
$\bigcap_{n=1}^{\infty}\bigl(-\tfrac{1}{n}, \tfrac{1}{n}\bigr) = \{0\}$
açık değildir.
\end{theorem}

\begin{theorem}[Baz Teoremi]
\label{thm:basis}
$\mathcal{B} \subseteq \mathcal{P}(X)$ ailesi \emph{(B1)--(B2)} koşullarını sağlıyorsa,
\[
  \tau_{\mathcal{B}}
  = \bigl\{U \subseteq X : \forall x \in U,\, \exists B \in \mathcal{B},\,
    x \in B \subseteq U\bigr\}
\]
bir topoloji olur ve $\mathcal{B}$ bu topolojiye baz oluşturur.
\end{theorem}

\begin{proof}[Kanıt İskeleti]
(T1): $\emptyset \in \tau_{\mathcal{B}}$ boş koşul nedeniyle; $X \in \tau_{\mathcal{B}}$
(B1) nedeniyle. (T2): $U, V \in \tau_{\mathcal{B}}$ ve $x \in U \cap V$ için
$B_U, B_V$ bul; (B2) ile $B_3 \subseteq B_U \cap B_V$. (T3): birleşim noktası için
doğrudan baz elemanı al. $\square$
\end{proof}

\begin{theorem}[Karşılaştırma]
\label{thm:comparison}
$X$ üzerinde iki topoloji $\tau_1 \subseteq \tau_2$ ise $\tau_1$ \textbf{kaba (coarser)},
$\tau_2$ \textbf{ince (finer)} topoloji olarak adlandırılır. Ayrık topoloji en ince,
indirgenmiş topoloji en kaba topolojidir.
\end{theorem}

%% ============================================================
\section{Algoritmalar}
%% ============================================================

\subsection{Bazdan Topoloji Üretimi}

\textbf{Problem:} $X$ ve $\mathcal{B} \subseteq \mathcal{P}(X)$ verilsin;
$\tau_{\mathcal{B}}$'yi hesapla.

\begin{algorithm}
\caption{Bazdan Topoloji Üretimi}
\label{alg:basis_to_topology}
\begin{algorithmic}[1]
\Require $X$ (taşıyıcı), $\mathcal{B}$ (baz ailesi)
\Ensure $\tau$ (kapalı topoloji)
\State $\tau \leftarrow \{\emptyset, X\}$
\For{each non-empty $\mathcal{S} \subseteq \mathcal{B}$}
    \State $\tau.\text{add}\bigl(\bigcup_{B \in \mathcal{S}} B\bigr)$
\EndFor
\State \Call{CloseUnderIntersection}{$\tau$}
\Return $\tau$
\end{algorithmic}
\end{algorithm}

\begin{algorithm}
\caption{Kesişim Altında Kapatma}
\begin{algorithmic}[1]
\State $\textit{changed} \leftarrow \textbf{true}$
\While{$\textit{changed}$}
    \State $\textit{changed} \leftarrow \textbf{false}$
    \For{each $U, V \in \tau$}
        \If{$U \cap V \notin \tau$}
            \State $\tau.\text{add}(U \cap V)$;\quad $\textit{changed} \leftarrow \textbf{true}$
        \EndIf
    \EndFor
\EndWhile
\end{algorithmic}
\end{algorithm}

\noindent\textbf{Karmaşıklık:}
\begin{itemize}
  \item En kötü durum: $O(|\mathcal{B}| \cdot 2^{|\mathcal{B}|})$
        (tüm alt-ailelerin birleşimi)
  \item Pratik (çakışmayan baz elemanları): $O(|\tau| \cdot |\mathcal{B}|)$
  \item Uzay: $O(|\tau|)$ açık küme saklar
\end{itemize}

\subsection{Alt-Bazdan Topoloji Üretimi}

$\mathcal{S}$ verilsin; önce $\mathcal{S}$'nin sonlu kesişimlerinden baz oluştur,
ardından Algoritma~\ref{alg:basis_to_topology}'yi uygula.

\noindent\textbf{Karmaşıklık:} $O(|\mathcal{S}|^k \cdot 2^{|\mathcal{S}|^k})$ —
$k$: kesişim derinliği.

%% ============================================================
\section{pytop API}
%% ============================================================

\subsection{Sınıflar}

\begin{lstlisting}[language=Python, caption={Temel sınıf importları}]
from pytop import TopologicalSpace, FiniteTopologicalSpace
\end{lstlisting}

\begin{table}[h]
\centering
\begin{tabular}{lll}
\toprule
\textbf{Sınıf} & \textbf{Taşıyıcı} & \textbf{Topoloji} \\
\midrule
\texttt{TopologicalSpace} & \texttt{Any} & \texttt{Any} (None olabilir) \\
\texttt{FiniteTopologicalSpace} & sonlu \texttt{frozenset} & \texttt{frozenset[frozenset]} \\
\bottomrule
\end{tabular}
\caption{Temel uzay sınıfları}
\end{table}

\noindent\textbf{Alanlar:} \texttt{carrier} (altta yatan küme), \texttt{topology}
(açık kümeler), \texttt{tags} (özellik etiketleri), \texttt{metadata} (ek bilgiler).

\subsection{Oluşturucular}

\begin{lstlisting}[language=Python, caption={Topoloji oluşturucu importları}]
from pytop import (
    make_topology, discrete_topology, indiscrete_topology,
    cofinite_topology, sierpinski_space,
    topology_from_basis, topology_from_subbasis,
)
\end{lstlisting}

\begin{table}[h]
\centering
\begin{tabular}{lll}
\toprule
\textbf{Fonksiyon} & \textbf{İmza} & \textbf{Açıklama} \\
\midrule
\texttt{make\_topology} & \texttt{(carrier, *open\_sets)} & Açık küme listesiyle inşa \\
\texttt{discrete\_topology} & \texttt{(*elements)} & Ayrık topoloji \\
\texttt{indiscrete\_topology} & \texttt{(*elements)} & İndirgenmiş topoloji \\
\texttt{cofinite\_topology} & \texttt{(*elements)} & Kosonlu topoloji \\
\texttt{sierpinski\_space} & \texttt{()} & Sierpiński uzayı \\
\texttt{topology\_from\_basis} & \texttt{(carrier, basis)} & Bazdan topoloji üret \\
\texttt{topology\_from\_subbasis} & \texttt{(carrier, subbasis)} & Alt-bazdan topoloji üret \\
\bottomrule
\end{tabular}
\caption{Topoloji oluşturucu fonksiyonlar}
\end{table}

\subsection{Tag Sistemi}

\texttt{pytop} nesneleri \texttt{tags} kümesi aracılığıyla topolojik özellikler taşır.
Örnek etiketler: \texttt{"finite"}, \texttt{"discrete"}, \texttt{"t0"}, \texttt{"t1"},
\texttt{"hausdorff"}, \texttt{"compact"}, \texttt{"connected"}, \texttt{"metric"}.

%% ============================================================
\section{Örnekler}
%% ============================================================

\subsection{Örnek 1: Sierpiński Uzayı}

\begin{lstlisting}[language=Python, caption={Sierpiński uzayı}]
from pytop import sierpinski_space

s = sierpinski_space()
print("Tasiyici:", s.carrier)
print("Topoloji:", sorted(str(t) for t in s.topology))
print("Etiketler:", sorted(s.tags))
\end{lstlisting}

\begin{lstlisting}[caption={Çıktı}]
Tasiyici: frozenset({0, 1})
Topoloji: ['frozenset()', 'frozenset({0, 1})', 'frozenset({1})']
Etiketler: ['compact', 'connected', 'finite', 't0']
\end{lstlisting}

Topoloji tam olarak üç açık küme içerir: $\emptyset$, $\{1\}$, $\{0,1\}$.
$\{0\}$ açık değildir — bu T0 fakat T1 olmayan özelliğin kaynağıdır.

\subsection{Örnek 2: Ayrık Topoloji}

\begin{lstlisting}[language=Python, caption={Ayrık topoloji}]
from pytop import discrete_topology

d = discrete_topology(1, 2, 3)
print("|tau| =", len(d.topology))
print("Etiketler:", sorted(d.tags))
\end{lstlisting}

\begin{lstlisting}[caption={Çıktı}]
|tau| = 8
Etiketler: ['discrete', 'finite', 'hausdorff', 'metrizable', 'normal', 'regular']
\end{lstlisting}

$n=3$ için $|\tau| = 2^3 = 8$. Ayrık topoloji Hausdorff, regüler, normal, metriklenebilirdir.

\subsection{Örnek 3: \texttt{make\_topology} ile Manuel İnşa}

\begin{lstlisting}[language=Python, caption={Manuel topoloji inşası}]
from pytop import make_topology

sp = make_topology({1, 2, 3}, {1}, {2, 3})
print("|tau| =", len(sp.topology))
print("Acik kumeler:", sorted(str(t) for t in sp.topology))
\end{lstlisting}

\begin{lstlisting}[caption={Çıktı}]
|tau| = 4
Acik kumeler: ['frozenset()', 'frozenset({1, 2, 3})',
               'frozenset({1})', 'frozenset({2, 3})']
\end{lstlisting}

$\{1\}$ ve $\{2,3\}$ açık verilince, $\emptyset$ ve $X$ otomatik eklenir.

\subsection{Örnek 4: Alexandrov Zincir Uzayı}

\begin{lstlisting}[language=Python, caption={Zincir uzayı}]
from pytop import finite_chain_space

c = finite_chain_space(3)
print("Tasiyici:", c.carrier)
print("Topoloji:", sorted(str(t) for t in c.topology))
\end{lstlisting}

\begin{lstlisting}[caption={Çıktı}]
Tasiyici: (1, 2, 3)
Topoloji: ['set()', '{1, 2, 3}', '{1, 2}', '{1}']
\end{lstlisting}

Her açık küme bir ``önek''tir. Alexandrov topolojisinde noktalar soldan sağa giderek az açık hale gelir.

\subsection{Örnek 5: Bazdan Topoloji Üretimi}

\begin{lstlisting}[language=Python, caption={Bazdan topoloji}]
from pytop import topology_from_basis

b = [{1}, {2, 3}, {4}]
ts = topology_from_basis({1, 2, 3, 4}, b)
print("Baz:", [set(x) for x in b])
print("|tau| =", len(ts.topology))
print("Topoloji:", sorted(str(t) for t in ts.topology))
\end{lstlisting}

\begin{lstlisting}[caption={Çıktı}]
Baz: [{1}, {2, 3}, {4}]
|tau| = 8
Topoloji: ['set()', '{1, 2, 3, 4}', '{1, 2, 3}', '{1, 4}',
           '{1}', '{2, 3, 4}', '{2, 3}', '{4}']
\end{lstlisting}

Baz $\{\{1\}, \{2,3\}, \{4\}\}$ bir bölüntüden oluşur; 8 açık küme tüm birleşim kombinasyonlarından elde edilir.

\subsection{Örnek 6: Gerçek Doğru}

\begin{lstlisting}[language=Python, caption={Gerçek doğru}]
from pytop import real_line_metric

rl = real_line_metric()
print("Tasiyici:", rl.carrier)
print("Etiketler:", sorted(rl.tags))
\end{lstlisting}

\begin{lstlisting}[caption={Çıktı}]
Tasiyici: R
Etiketler: ['complete', 'connected', 'first_countable', 'hausdorff',
            'infinite', 'lindelof', 'metric', 'not_compact',
            'path_connected', 'second_countable', 'separable', 't0', 't1', 'uncountable']
\end{lstlisting}

Gerçek doğru simgesel temsil edilir; topolojik özellikler etiketlerde kodlanmıştır.

%% ============================================================
\section{Alıştırmalar}
%% ============================================================

\subsection*{Kodlama Alıştırmaları}

\begin{enumerate}[label=\textbf{K\arabic*.}]
  \item $X = \{a, b, c\}$ üzerinde \texttt{make\_topology} kullanarak T1 ama T2 (Hausdorff)
        olmayan bir topoloji inşa edin. (İpucu: kosonlu topolojiyi düşünün.)

  \item \texttt{topology\_from\_subbasis(\{1,2,3,4\}, [\{1,2\},\{3,4\},\{2,3\}])} çağrısının
        ürettiği topolojinin kaç açık küme içerdiğini bulun ve açık kümeleri listeleyin.

  \item \texttt{finite\_chain\_space(5)} ile bir Alexandrov-5 zinciri oluşturun; topolojinin
        kaç açık küme içerdiğini ve hangi elemanların ``en açık'' nokta olduğunu bulun.
\end{enumerate}

\subsection*{Teori Alıştırmaları}

\begin{enumerate}[label=\textbf{T\arabic*.}]
  \item $X = \{1, 2, 3\}$ üzerinde kaç farklı topoloji tanımlanabilir? Bu sayıyı T1--T3
        aksiyomlarını doğrudan kontrol ederek hesaplamak yerine neden baz teoremini kullanmak
        daha pratiktir?

  \item Ayrık topolojinin her zaman diğer topolojilerden daha ince, indirgenmiş topolojinin
        ise daha kaba olduğunu kanıtlayın. Kanıt için T1--T3 aksiyomlarından hangisi gereklidir?
\end{enumerate}
